\chapter{The interpolation and the defect}\label{sec:fnd}

In this section, we recall the interpolation of \cite[Section~2]{AL},
which compares the fractional moments of the partition function with the
constrained scalar Parisi problem. We also fix the normalizations, the
tilted law, and the observables that we use in the rest of the
paper. The later sections build on two results. The first is the
defect identity \eqref{eq:in-envelope} of \cref{cor:fnd-defect}. It states
that four times the fractional-moment error $e_{N,\beta}(\lambda)$
is the integral over the interpolation of a defect
$\calD_N(\rho)=\beta(A+2J+V)$, whose three terms are overlap observables
at a minimizing profile. The second is \cref{prop:fnd-mean-bias}, where
we show that $(P_\beta-p_{N,\beta})/\beta$ is at
most $2e_{N,\beta}(\lambda)$ plus a term of order $\lambda^5$. Along the
way, we record the martingale and contact identities in
\cref{prop:fnd-contact}, the integrated variance bound in
\cref{prop:fnd-variance}, the decomposition of the overlap error in
\cref{lem:fnd-decomp}, the interpolation identity in
\cref{prop:fnd-interpolation}, and the trial budget in
\cref{lem:fnd-trial}.

Most of these facts are proved in \cite{AL}. We state each of them in
full, with its location in \cite{AL}. Where the argument is short, we
also sketch the proof, so that this section can be read without \cite{AL}
at hand. Besides \cite{AL}, this section uses only the scalar Parisi
functional \eqref{eq:intro-parisi} and the notation of
\cref{sec:intro-background}, together with the support theorem
\cite[Theorem~1.1]{Lopatto}. Its results are used throughout the
paper. \Cref{prop:fnd-contact,prop:fnd-variance,lem:fnd-decomp} are
used in \cref{sec:wi,sec:th,sec:c7,sec:c7-closure}; the energy
normalization and the trial budget in \cref{sec:lz,sec:lz-root,sec:sc};
and \eqref{eq:in-envelope}, \eqref{eq:fnd-moving}, and
\cref{prop:fnd-mean-bias} in \cref{sec:c7-closure,sec:b7}. We treat the
laws with anisotropic fields in \cref{sec:gauss}, where we
also prove the two Gaussian tools that we need for profiles that do not
minimize the functional.

Throughout this section, we fix $\beta>1$ and $N\ge1$, and we assume no
relation between the tilt $\lambda$ and $N$. The completion,
the interpolating functional, and the identities of
\cref{subsec:fnd-law,sec:fnd-contact,sec:fnd-variance,sec:fnd-defect}
hold for every $\beta>0$. We use the hypothesis $\beta>1$ for the support
of the Parisi measure in \cref{sec:fnd-completion} and for the
cost of the scalar floor in \cref{sec:fnd-trial}.

The interpolation of \cite[Section~2]{AL} is a quantitative form of the
Guerra interpolation that Talagrand used for the positive fractional
moments of the partition function, whose limits he determined in
\cite{Talagrand2007FractionalMoments}; see \cite[Section~1.2]{AL}. For
fixed $0<s<1$, the limit of the normalized logarithmic moment of order
$s$ is the minimum of the Parisi functional over probability measures
with mass at least $s$ at zero. For the SK model at zero field, this
formula is also contained in the one-level case of
\cite[Theorem~2.1]{ContucciMingione2019}, and it is stated in this form
in \cite[Theorem~2.1]{ChenEtAl2026LargeDeviations}, in a normalization
of the spins that changes the functional by the constant $\log2$. The
finite-$N$ functional of \cref{sec:fnd-functional} carries the same
constraint, in the sense that its profiles have the
prescribed mass $s$ at $0$.

\section{Completion and the constrained scalar value}
\label{sec:fnd-completion}

Let $W$ be the symmetric $N\times N$ matrix with entries
$W_{ij}=g_{ij}/\sqrt N$ for $i<j$ and $W_{ii}=\sqrt{2/N}\,z_i$, where
$z_1,\ldots,z_N$ are standard Gaussian variables, independent of each
other and of the couplings $(g_{ij})$. We write $\E_{\rm ref}$ for
expectation over $(g_{ij},z_i)$ under this Gaussian law; the subscript
distinguishes it from the tilted law of \cref{subsec:fnd-law}.
Following \cite[Section~1.1]{AL}, we set
\[
 D_N=\frac12\Tr W,\qquad
 F^c_{N,\beta}=\log\sum_{\sigma\in\{-1,1\}^N}e^{\beta\sigma^TW\sigma/2},
\]
and we recall that $F_{N,\beta}=\log\sum_\sigma e^{\beta H_N(\sigma)}$ is the
logarithmic partition function of \cref{sec:intro-background}. Note that
$\sigma^TW\sigma/2=H_N(\sigma)+D_N$, because $\sigma_i^2=1$, so that
$F^c_{N,\beta}=F_{N,\beta}+\beta D_N$. The variable
$D_N=(2N)^{-1/2}\sum_iz_i$ is centered Gaussian with variance $1/2$ and
is independent of $F_{N,\beta}$.

The diagonal term completes the covariance of the Hamiltonian to an
exact function of the overlap $R(\sigma,\sigma')=N^{-1}\sigma\cdot\sigma'$.
Indeed, writing $\tau_i=\sigma_i\sigma'_i$, we have
\[
 \E_{\rm ref}\bigl[H_N(\sigma)H_N(\sigma')\bigr]
 =\frac1N\sum_{i<j}\tau_i\tau_j
 =\frac1{2N}\biggl[\Bigl(\sum_{i=1}^N\tau_i\Bigr)^2-N\biggr]
 =\frac N2R(\sigma,\sigma')^2-\frac12 ,
\]
and adding $\Var(D_N)=1/2$, we find that the completed Hamiltonian
$\sigma^TW\sigma/2$ has covariance $NR(\sigma,\sigma')^2/2$. This is why
the derivative of the interpolation in \cref{prop:fnd-interpolation} is
exactly an integral of squared overlap deviations. Because $D_N$ is
independent of $F_{N,\beta}$, the completion changes the logarithmic
moments only by an explicit term. For $s>0$, the normalized logarithmic moment of the completed
partition function is \cite[(1.3)]{AL}
\[
 \phi^c_{N,\beta}(s)
   =\frac1{Ns}\log\E_{\rm ref}e^{sF^c_{N,\beta}}
   =\frac1{Ns}\log\E_{\rm ref}e^{sF_{N,\beta}}+\frac{\beta^2s}{4N},
\]
where we used $\log\E_{\rm ref}e^{s\beta D_N}=\beta^2s^2/4$. As $s\downarrow0$, we have $\phi^c_{N,\beta}(s)\to p_{N,\beta}$, since
$F_{N,\beta}$ has exponential moments of every order.

Recall the scalar Parisi functional $\mathcal P_\beta$ of
\eqref{eq:intro-parisi}, which we evaluate at the distribution function
of a probability measure on $[0,1]$. For $0\le s\le1$ and
$0\le\lambda\le\beta$, we define
\[
 P_\beta(s)=\min_{\mu\in\mathcal P([0,1]):\,\mu(\{0\})\ge s}
      \mathcal P_\beta\bigl(\mu([0,\cdot\,])\bigr),
 \qquad E_\beta(\lambda)=\beta^{-1}P_\beta(\lambda/\beta).
\]
The minimum is attained \cite[Section~1.1]{AL}, and $P_\beta(0)=P_\beta$.
The function $E_\beta$ is the constrained value in the energy
normalization of \cite[Section~7.1]{AL}, in which the order $s$ of the
moment is replaced by the tilt $\lambda=\beta s$. Note that $E_\beta$ is nondecreasing, because the admissible class
shrinks as $s$ grows; in particular, $E_\beta(\lambda)\ge E_\beta(0)$.

The unconstrained minimizer $\mu_\beta$ is unique
\cite{AuffingerChen2015Uniqueness}. We write
$\alpha_\beta(q)=\mu_\beta([0,q])$ and $\gamma_\beta=\beta\alpha_\beta$.
Since $\beta>1$, we have $\supp\mu_\beta=[0,q_\beta]$ with
$0<q_\beta<1$, by \cite[Theorem~1.1]{Lopatto}. We recall the further
structure of $\mu_\beta$ in \cref{sec:lz}.

For $0<\lambda<\beta$ and $s=\lambda/\beta$, we define the
\emph{fractional-moment error} by
\begin{equation}\label{eq:fnd-error-definition}
 e_{N,\beta}(\lambda)
    =\frac{P_\beta(s)-\phi^c_{N,\beta}(s)}\beta
    =E_\beta(\lambda)
       -\frac1{N\lambda}\log\E_{\rm ref}
          e^{\lambda(F_{N,\beta}/\beta+D_N)} .
\end{equation}
In the second expression, we used
$sF^c_{N,\beta}=\lambda(F_{N,\beta}/\beta+D_N)$; the completion
contributes $\lambda/(4N)$ to the logarithmic moment on the right.
As $\beta$ is fixed, we also write $e_N(\lambda)=e_{N,\beta}(\lambda)$.
From now on, we assume that $s=\lambda/\beta\le1/2$, which is the range
of the results of \cite[Section~2]{AL}. We will see in
\cref{cor:fnd-defect} that the error
$e_{N,\beta}(\lambda)$ is nonnegative.

\section{The interpolating functional}\label{sec:fnd-functional}

For $0\le\rho\le1$, we set
\[
 t=1-\rho,\qquad c=\beta^2\rho,\qquad M=Nc .
\]
The interpolation reaches the fractional moment at $\rho=0$, and
$\rho=1$ is the scalar end. We parametrize the interpolation by $\rho$
because the estimates of the later sections degenerate as
$\rho\downarrow0$. We work with the tilted law of \cref{subsec:fnd-law}
for $0<\rho\le1$,
where $c>0$.

An \emph{admissible profile} is a deterministic probability measure
$\mu=s\delta_0+\nu$ on $[0,1]$, where $\nu\ge0$ has mass $1-s$ and may
charge $0$. Equivalently, $\mu$ is a deterministic probability measure
such that $\mu(\{0\})\ge s$. We call $s\delta_0$ the prescribed root mass and
$\nu$ the free measure of $\mu$. We write $a(q)=\mu([0,q])$ for the
cumulative function. It is
nondecreasing and right-continuous, with $a\ge s$ and $a(1)=1$.
Conditional on $W$, let $U=U_{\mu,t}$ solve
\begin{equation}\label{eq:fnd-vector-pde}
 \partial_qU=-\frac c2\bigl(\Delta U+a(q)|\nabla U|^2\bigr),
 \qquad
 U(1,x;W)=\log\sum_{\sigma\in\{-1,1\}^N}
  e^{\beta\sqrt t\,\sigma^TW\sigma/2+x\cdot\sigma}.
\end{equation}
This is the Parisi equation in $\R^N$; its terminal value is the logarithm
of the completed partition function at inverse temperature $\beta\sqrt t$
in the external field $x$. We define the interpolating functional and its
minimum over admissible
profiles by
\begin{equation}\label{eq:fnd-functional}
 \mathcal F_{t,s}(\mu)
   =\frac1{Ns}\log\E_{\rm ref}e^{sU(0,0;W)}
       -\frac c2\int_0^1qa(q)\dd q,
 \qquad f_{N,s}(t)=\min_{\mu}\mathcal F_{t,s}(\mu).
\end{equation}

By \cite[Section~2.6 and Proposition~2.3]{AL}, the equation
\eqref{eq:fnd-vector-pde} has a solution for every admissible profile,
and $\mathcal F_{t,s}$ is continuous under weak convergence of profiles.
For a step profile, the solution is given by the Hopf--Cole recursion that
we recall in
\cref{subsec:fnd-law}. For a general profile, it is the
uniform limit of the solutions for step profiles whose cumulative
functions converge in $L^1([0,1])$. Indeed, if $U_a$ and $U_{\tilde a}$
are the solutions for the cumulative functions $a$ and $\tilde a$, then
\cite[Section~2.6, (2.19)]{AL}
\[
 \sup_{q,x,W}\bigl|U_a(q,x;W)-U_{\tilde a}(q,x;W)\bigr|
 \le\frac{cN}2\int_0^1|a(u)-\tilde a(u)|\dd u .
\]
The admissible profiles form a weakly compact set, since they are the
image of the compact set of measures of mass $1-s$ under
$\nu\mapsto s\delta_0+\nu$, and weak convergence of profiles implies
$L^1$ convergence of their cumulative functions. We conclude that the
minimum in \eqref{eq:fnd-functional} is attained.

We note that minimizers are deterministic. The Gaussian expectation is
taken inside
the functional, which means that $\mathcal F_{t,s}(\mu)$ is a number for
each profile. A minimizer may depend on $N$, $\rho$, and $\lambda$ but not
on the couplings. This matters below, because the tilted law of
\cref{subsec:fnd-law}
is constructed from a deterministic profile, and the estimates at a
minimizer combine that law with the variational inequality that the same
profile satisfies. A minimizer need not be unique; at each use, we say
which minimizer is meant.

The prescribed root mass $s\delta_0$ matches the order $s$ of the
moment. In the derivative of the interpolation, which we compute in
\cref{prop:fnd-interpolation}, the outer logarithmic moment
contributes a term that cancels the root mass exactly, so that only the
free measure $\nu$ remains. In the posterior densities of
\cref{sec:gauss}, the outer tilt cancels it in the same way.

The two endpoints of the interpolation are explicit. At $t=0$, the
terminal value in \eqref{eq:fnd-vector-pde} is $\sum_i\log(2\cosh x_i)$
and $c=\beta^2$. Then the equation separates into $N$ copies of the
scalar Parisi equation of \eqref{eq:intro-parisi}, and we have
$U(q,x;W)=\sum_iu_a(q,x_i)$, where $u_a$ is the solution of that
equation with $\alpha=a$. This shows that $U(0,0;W)=Nu_a(0,0)$ is
deterministic and that
$\mathcal F_{0,s}(\mu)=\mathcal P_\beta(a)$. At
$t=1$, we have $c=0$, which implies that
$U(0,0;W)=U(1,0;W)=F^c_{N,\beta}$ and that
$\mathcal F_{1,s}(\mu)=\phi^c_{N,\beta}(s)$ for every profile.
Minimizing over admissible profiles, we obtain
\cite[Proposition~2.3]{AL}
\begin{equation}\label{eq:fnd-endpoints}
 f_{N,s}(0)=P_\beta(s),\qquad f_{N,s}(1)=\phi^c_{N,\beta}(s).
\end{equation}
For a fixed profile, the map $t\mapsto\mathcal F_{t,s}(\mu)$ is
nonincreasing (\cref{prop:fnd-interpolation}). By
\eqref{eq:fnd-error-definition}, the fractional-moment error is
$(f_{N,s}(0)-f_{N,s}(1))/\beta$, which is the decrease of $f_{N,s}/\beta$
along the interpolation. The defect of \cref{sec:fnd-defect} is four
times the rate of this decrease.

\paragraph{The energy normalization.}
We vary profiles in the energy normalization of
\cite[Section~7.1]{AL}. For an admissible profile $\mu$ with cumulative
function $a$, we set
\[
 \gamma=\beta a,\qquad \nu_\gamma=\beta\nu,\qquad
 F_\rho(\gamma)=\beta^{-1}\mathcal F_{1-\rho,\lambda/\beta}(\mu).
\]
Then $\dd\gamma$ is a measure of mass $\beta$ with
$\dd\gamma\ge\lambda\delta_0$; we also call $\gamma$ admissible. We
have $\min_\gamma F_\rho(\gamma)=\beta^{-1}f_{N,s}(1-\rho)$, where the
minimum is over admissible $\gamma$. The name refers to the following
rescaling of the spatial variable. The function
$\widetilde U(q,y;W)=\beta^{-1}U(q,\beta y;W)$ solves
\[
 \partial_q\widetilde U
 =-\frac\rho2\bigl(\Delta\widetilde U+\gamma(q)|\nabla\widetilde U|^2\bigr),
\]
with terminal value $\beta^{-1}$ times the logarithm of the completed
partition function in the field $\beta y$, and we have
\[
 F_\rho(\gamma)=\frac1{N\lambda}\log\E_{\rm ref}
     e^{\lambda\widetilde U(0,0;W)}
 -\frac\rho2\int_0^1q\gamma(q)\dd q .
\]
This is the functional of \cite[(7.2)]{AL} with $r=\rho$, in which the
diagonal of $W$ has been integrated out explicitly
\cite[Section~7.1]{AL}. Note that the rescaling leaves the gradient unchanged,
$\nabla\widetilde U(q,y;W)=\nabla U(q,\beta y;W)$, and that it multiplies
the Hessian by $\beta$.

We define the \emph{scalar trial profile} by
$\gamma^{\rm tr}_\lambda(q)=\max\{\lambda,\gamma_\beta(q)\}$. It is
nondecreasing and right-continuous, and it satisfies
$\lambda\le\gamma^{\rm tr}_\lambda\le\beta$ and
$\gamma^{\rm tr}_\lambda(1)=\beta$. In other words,
$\gamma^{\rm tr}_\lambda/\beta$ is the cumulative function of an
admissible profile. When $\lambda$ is
fixed, we write $\gamma^{\rm tr}=\gamma^{\rm tr}_\lambda$. We discuss its
cost and its use as a comparison profile in
\cref{sec:fnd-trial}.

\section{The tilted joint law and overlap observables}
\label{subsec:fnd-law}

Fix $0<\rho\le1$ and an admissible profile $\mu$. We recall the joint
law of \cite[Section~2]{AL}, under which all later expectations are
taken. It is built in three stages.
\begin{enumerate}[label=(\roman*)]
\item The Gaussian law of $W$ is tilted by the density
$e^{sU(0,0;W)}/\E_{\rm ref}e^{sU(0,0;W)}$.
\item Given $W$, the auxiliary field solves
\begin{equation}\label{eq:fnd-diffusion}
 \dd X_q=ca(q)\nabla U(q,X_q;W)\dd q+\sqrt c\dd B_q,
 \qquad X_0=0,
\end{equation}
where $B$ is a standard Brownian motion in $\R^N$ under the conditional
law given $W$.
\item Given $W$ and the path $X$, a terminal spin
$\sigma\in\{-1,1\}^N$ is drawn with probabilities proportional to
$\exp\{\beta\sqrt t\,\sigma^TW\sigma/2+X_1\cdot\sigma\}$.
\end{enumerate}
We write $\Pp$ and $\E$ for this joint law, and
$\calG_q=\sigma(W,X_r:0\le r\le q)$ for the information revealed up to
level $q$; note that it does not contain the terminal spin. At fixed $N$, the
drift in \eqref{eq:fnd-diffusion} is bounded and Lipschitz in $x$
\cite[Section~2.6.1]{AL}, and the equation has a unique strong solution.
The tilt in (i) changes the law of $W$, which is no longer Gaussian under
$\Pp$, but it does not change the
conditional law of $B$ given $W$.

For a step profile, the law has a discrete description, which explains
the construction. The Gaussian Hopf--Cole recursion in the Parisi setting
is described in \cite[Section~3, (10)--(13)]{Guerra2003}, and the
finite-dimensional tilted law used here is the formulation of
\cite[Section~2.1]{AL}.
Suppose that the profile has grid points
$0=q_0<q_1<\cdots<q_K=1$ and that $a=a_\ell$ on $[q_\ell,q_{\ell+1})$. Then
\eqref{eq:fnd-vector-pde} is solved by the Hopf--Cole recursion
\[
 U(q_\ell,x;W)=\frac1{a_\ell}\log\E_g
 \exp\bigl\{a_\ell U\bigl(q_{\ell+1},x+\sqrt{c(q_{\ell+1}-q_\ell)}\,g;W\bigr)\bigr\},
\]
where $g$ is a standard Gaussian vector in $\R^N$, and we read the right
side as $\E_gU(q_{\ell+1},x+\sqrt{c(q_{\ell+1}-q_\ell)}\,g;W)$ when
$a_\ell=0$. Given $W$ and $X_{q_\ell}$, the increment
$X_{q_{\ell+1}}-X_{q_\ell}$ has the law of
$\sqrt{c(q_{\ell+1}-q_\ell)}\,g$, tilted by the density
$\exp\{a_\ell(U(q_{\ell+1},X_{q_{\ell+1}};W)-U(q_\ell,X_{q_\ell};W))\}$,
which integrates to one by the recursion. The outer tilt in (i) has the
same form, with mass $s$ at level $0$. In \cref{sec:gauss}, we multiply
these densities and use the resulting formula for the joint density.

\paragraph{Observables.}
We denote the conditional mean, the conditional covariance, and the
spatial Hessian by
\[
 m_q=\E[\sigma\mid\calG_q],\qquad
 C_q=\Cov(\sigma\mid\calG_q),\qquad
 H_q=\nabla_x^2U(q,X_q;W).
\]
We will see in \cref{prop:fnd-contact} that $m_q=\nabla U(q,X_q;W)$ and
that $a(q)C_q\preceq H_q\preceq C_q$. With the same terminal spin at every
level, we define
\[
 S_q=\frac{|m_q|^2}N,\qquad
 L_q=\frac{\sigma\cdot m_q-|m_q|^2/2}N,\qquad
 j_q=\frac{\E[m_q^TC_qm_q]}{N^2},\qquad
 a_q^{\rm th}=\frac{\E\Tr C_q^2}{N^2};
\]
we use the superscript in $a^{\rm th}_q$ to distinguish the covariance
trace from the cumulative function $a$. For each $q$, let $\sigma^1$ and $\sigma^2$
be conditionally independent given $\calG_q$, each with the conditional
law of $\sigma$ given $\calG_q$, and let $R_q=R(\sigma^1,\sigma^2)$. Note
that the self-overlap $S_q$ of the conditional mean is the conditional
expectation of $R_q$ given $\calG_q$, and that $L_q$ has conditional
expectation $S_q/2$ given $\calG_q$. In \cref{prop:fnd-variance}, we bound
the variance of
averages of $L_q$ over levels. We also have
$0\le S_q\le1$ and $|L_q|\le3/2$.

The processes $m$, $C$, and $H$ have continuous paths. Indeed, we have
$m_q=\nabla U(q,X_q;W)$, $H_q=\nabla^2U(q,X_q;W)$, and
$\E[\sigma\sigma^T\mid\calG_q]=\Sigma(q,X_q;W)$, where the conditional
second moment $\Sigma$ solves a linear backward equation with drift
$ca\nabla U$, and $\nabla U$, $\nabla^2U$, and $\Sigma$ are jointly
continuous in $(q,x)$ \cite[Section~2.6.1]{AL}; the path $X$ is
continuous. In particular, $S_q$ and $L_q$ are jointly measurable
in the level and the outcome, and Fubini's theorem applies to their
integrals against $\nu$. By bounded convergence, the functions
$q\mapsto j_q$, $a^{\rm th}_q$, $\E S_q$, $\Var(S_q)$, and $\E|S_q-q|^3$
are continuous on $[0,1]$, and so is $q\mapsto\E(R_q-q)^2$, because
$\E[R_q\mid\calG_q]=S_q$ and
$\E[R_q^2\mid\calG_q]=N^{-2}\sum_{i,j}\Sigma_{ij}(q,X_q;W)^2$.

We also define the integrated quantities
\[
 \begin{gathered}
 A=\int a_q^{\rm th}\,\nu(\dd q),\qquad
 J=\int j_q\,\nu(\dd q),\\
 V=\int\Var(S_q)\,\nu(\dd q),\qquad
 H_3=\int\E|S_q-q|^3\,\nu(\dd q).
 \end{gathered}
\]
We note that they depend on $N$, $\rho$, $\lambda$, and the profile. The trace term
$A$ and the mean--covariance term $J$ are controlled by the conditional
covariance. The integrated variance $V$ of the self-overlap is the main
term of the defect, and $H_3$ is the corresponding cubic moment, which
enters in \cref{sec:c7-closure}.

\paragraph{Quantiles and contact levels.}
The quantile function
\begin{equation}\label{eq:fnd-quantile}
 Q(v)=\inf\{q\in[0,1]:\nu([0,q])\ge v\},\qquad 0<v<1-s,
\end{equation}
pushes Lebesgue measure on $(0,1-s)$ forward to $\nu$, and
$a(Q(v))=s+\nu([0,Q(v)])\ge s+v$, by the right continuity of
$q\mapsto\nu([0,q])$. We use quantiles to integrate against $\nu$ in a
variable in which the lower bound on $a$ is explicit, whether or not
$\nu$ has a density. The image under $Q$ of Lebesgue measure on a Borel
subset of $(0,1-s)$ is a sub-measure of $\nu$, whose density with respect
to $\nu$ takes
values in $[0,1]$. For an interval of values of
$v$, this sub-measure may contain only part of an atom of $\nu$, and by
integration over such an interval, we mean integration against
this sub-measure.

If $\mu$ minimizes $\mathcal F_{t,s}$, we call a point of $\supp\nu$ a
\emph{contact level} of $\mu$. The contact levels are where the free
measure sits, and \cref{prop:fnd-contact} shows that the expected self-overlap
$\E S_q$ equals $q$ at each of them.

\section{Martingale and contact identities}\label{sec:fnd-contact}

In the next proposition, we collect the finite-dimensional martingale,
first-variation, and contact identities of \cite[Proposition~2.1]{AL}.
For the scalar Parisi functional, the stochastic-control representation
and the martingale identities are in
\cite[Theorem~3 and Lemma~2]{AuffingerChen2015Uniqueness}, with a dynamic
programming treatment in \cite{JagannathTobasco2016}. The first variation
and the optimality conditions on the support are in
\cite[Lemma~3.2 and Proposition~1.1]{JagannathTobasco2017PhaseDiagram}
and \cite[Theorem~2 and Proposition~1]{Chen2017Variational}.

\begin{proposition}[Martingale, covariance order, and contact]
\label{prop:fnd-contact}
Let $0<\rho\le1$, and let $\mu=s\delta_0+\nu$ be an admissible profile.
For every $q\in[0,1]$,
\begin{equation}\label{eq:in-martingale}
 m_q=\nabla U(q,X_q;W),\qquad m_0=0,\qquad
 \dd m_q=\sqrt c\,H_q\dd B_q,
 \qquad a(q)C_q\preceq H_q\preceq C_q ,
\end{equation}
and $q\mapsto\E S_q$ is continuously differentiable on $[0,1]$ with
derivative $c\,\E\Tr H_q^2/N$. If $\nu'\ge0$ has mass $1-s$, then the
right derivative of
$\theta\mapsto\mathcal F_{t,s}(s\delta_0+(1-\theta)\nu+\theta\nu')$ at
$\theta=0$ is
\begin{equation}\label{eq:fnd-first-variation}
 D\mathcal F_{t,s}(\mu)[\nu'-\nu]
    =\frac c2\int_0^1G(x)\,(\nu'-\nu)(\dd x),\qquad
 G(x)=\int_x^1(\E S_y-y)\dd y .
\end{equation}
If $\mu$ minimizes $\mathcal F_{t,s}$, then $\supp\nu$ is contained in
the set of minimum points of $G$ on $[0,1]$, and
\begin{equation}\label{eq:in-contact}
 \supp\nu\subset[0,1),\qquad
 \E S_q=q,\qquad \E\Tr H_q^2\le\frac Nc
 \qquad\text{for every contact level }q .
\end{equation}
\end{proposition}

\begin{remark}\label{rem:fnd-root}
A contact level is a point of the support of the free measure $\nu$. At
$q=0$, we have $\E S_0=0$ for every profile, since $m_0=0$. The bound
on $\E\Tr H_0^2$ in \eqref{eq:in-contact} holds when $0\in\supp\nu$;
the prescribed mass $s\delta_0$ alone does not give it, because the
minimization does not act on that mass.
\end{remark}

\begin{proof}[Proof of \cref{prop:fnd-contact}]
We recall the argument of \cite[Proposition~2.1 and Section~2.6]{AL}.
The proof has four steps. We first prove the martingale identities and
the covariance order for step profiles. We then derive the first
variation and deduce the support property and \eqref{eq:in-contact}
at a minimizer. Finally, we extend these statements to general profiles.

\paragraph{Step 1: martingale identities for step profiles.}
Let $\mu$ be a step profile. Then $U$ is smooth in $x$, every spatial
derivative of fixed order is bounded uniformly in $(q,x,W)$ at fixed
$N$, and $U$ is smooth in $q$ between consecutive grid points
\cite[Section~2.6.1]{AL}. Differentiating \eqref{eq:fnd-vector-pde} in
$x$ and applying It\^o's formula along \eqref{eq:fnd-diffusion} on each
interval between grid points, we obtain \cite[(2.5)]{AL}
\[
 \dd\,\nabla U(q,X_q;W)=\sqrt c\,H_q\dd B_q,\qquad
 \dd H_q=-ca(q)H_q^2\dd q+\sqrt c\,\nabla H_q\dd B_q ,
\]
where $\nabla H_q\dd B_q$ denotes the contraction of the third spatial
derivative of $U$ at $(q,X_q)$ with $\dd B_q$. The stochastic integrals are martingales, since their integrands are
bounded. At $q=1$, the
gradient $\nabla U(1,X_1;W)$ is the Gibbs mean of the terminal spin
given $(W,X_1)$, that is, $\E[\sigma\mid\calG_1]$. By the tower property,
we have $\nabla U(q,X_q;W)=\E[\sigma\mid\calG_q]=m_q$ for all $q$.
Replacing $\sigma$ by $-\sigma$ in the sum that defines $U(1,x;W)$, we
see that $U(1,\cdot\,;W)$ is even. The equation preserves this symmetry,
which implies
that $U(q,\cdot\,;W)$ is even and that $m_0=\nabla U(0,0;W)=0$.

For the covariance order, we first note that $C_1=H_1$, because the
Hessian of the logarithm of a partition function is the covariance of
its Gibbs measure. Applying the martingale isometry to $m$, and taking
conditional expectations in the equation for $H$, we find
\[
 C_q=\E[H_1\mid\calG_q]+c\int_q^1\E[H_u^2\mid\calG_q]\dd u,\qquad
 H_q=\E[H_1\mid\calG_q]+c\int_q^1a(u)\E[H_u^2\mid\calG_q]\dd u .
\]
Observe that every matrix in these expressions is positive
semidefinite, and that $a(q)\le a(u)\le1$ for $u\ge q$. Subtracting, we obtain
\[
 C_q-H_q=c\int_q^1\bigl(1-a(u)\bigr)\E[H_u^2\mid\calG_q]\dd u\succeq0
\]
and
\[
 H_q-a(q)C_q=\bigl(1-a(q)\bigr)\E[H_1\mid\calG_q]
 +c\int_q^1\bigl(a(u)-a(q)\bigr)\E[H_u^2\mid\calG_q]\dd u\succeq0,
\]
which proves $a(q)C_q\preceq H_q\preceq C_q$ \cite[(2.6)]{AL}. Finally,
It\^o's formula gives $\dd|m_q|^2=2m_q\cdot\dd m_q+c\Tr H_q^2\dd q$.
Taking expectations, we find that the derivative of $q\mapsto\E S_q$ is
$c\,\E\Tr H_q^2/N$ \cite[(2.7)]{AL}, which is continuous in $q$.

\paragraph{Step 2: the first variation.}
Let $\nu'\ge0$ have mass $1-s$, and let
$\delta a(q)=(\nu'-\nu)([0,q])$, so that the profile
$s\delta_0+(1-\theta)\nu+\theta\nu'$ has cumulative function
$a+\theta\,\delta a$. The derivative $\dot U$ of the solution in
$\theta$ at $\theta=0$ solves the linearized equation
\[
 \partial_q\dot U+\frac c2\Delta\dot U+ca(q)\nabla U\cdot\nabla\dot U
 =-\frac c2\,\delta a(q)|\nabla U|^2,\qquad \dot U(1,\cdot\,;W)=0 .
\]
Its generator is the generator of \eqref{eq:fnd-diffusion}. By the
Feynman--Kac formula, we have
\[
 \dot U(0,0;W)=\frac c2\,\E\Bigl[\int_0^1\delta a(q)|m_q|^2\dd q
 \Bigm|W\Bigr].
\]
When we differentiate the outer logarithmic moment in
\eqref{eq:fnd-functional}, this identity is averaged under the tilted law
of $W$, and the penalty
contributes $-(c/2)\int_0^1q\,\delta a(q)\dd q$. We obtain
\[
 D\mathcal F_{t,s}(\mu)[\nu'-\nu]
 =\frac c2\int_0^1(\E S_q-q)\,\delta a(q)\dd q .
\]
We write $\delta a(q)=\int\ind\{x\le q\}\,(\nu'-\nu)(\dd x)$ and
exchange the order of integration, which is allowed because the
integrand is bounded and $\nu'-\nu$ has finite total variation. This
gives \eqref{eq:fnd-first-variation}, which is \cite[(2.8)]{AL}. We note
that its derivation does not use minimality. In the energy
normalization, the
identity is stated for every admissible profile in
\cite[Section~7.4, (7.22)]{AL}.

\paragraph{Step 3: support and contact at a minimizer.}
Let $\mu$ minimize $\mathcal F_{t,s}$, and let $\nu'\ge0$ have mass
$1-s$. The profiles $s\delta_0+(1-\theta)\nu+\theta\nu'$ with
$0\le\theta\le1$ are admissible, and $\theta=0$ is a minimum point of
$\mathcal F_{t,s}$ along them. This implies that the right derivative in
\eqref{eq:fnd-first-variation} is nonnegative. Taking
$\nu'=(1-s)\delta_x$, we find
\[
 (1-s)G(x)\ge\int G\dd\nu\qquad\text{for every }x\in[0,1].
\]
The function $G$ is continuous on $[0,1]$. Choosing for $x$ a minimum
point of $G$, we see that $\int(G-\min_{[0,1]}G)\dd\nu\le0$, so that
$G=\min_{[0,1]}G$ holds $\nu$-almost everywhere. The set where
$G>\min_{[0,1]}G$ is relatively open in $[0,1]$ and $\nu$-null, and we
conclude that $\supp\nu$ is contained in the set of minimum points of
$G$.

By Step 1, $G$ is twice continuously differentiable, with
\[
 G'(q)=q-\E S_q,\qquad G''(q)=1-\frac cN\E\Tr H_q^2 .
\]
If $q\in\supp\nu\cap(0,1)$, then $q$ is an interior minimum point of
$G$, and therefore $G'(q)=0$ and $G''(q)\ge0$. These are the last two
relations in
\eqref{eq:in-contact}. At $q=0$, we have $G'(0)=-\E S_0=0$, because
$m_0=0$. If $0\in\supp\nu$, then $0$ is a minimum point of $G$, and
$G''(0)\ge0$, since otherwise $G''<0$ near $0$ and $G(h)<G(0)$ for
small $h>0$. At $q=1$, the terminal Gibbs measure gives positive
probability to every spin configuration. Then $|m_{1,i}|<1$ for all $i$,
and $G'(1)=1-\E S_1>0$, which implies that $G(1-h)<G(1)$ for small $h>0$
and that $1$ is not a minimum point of $G$. This proves
$\supp\nu\subset[0,1)$ and \eqref{eq:in-contact}, which is
\cite[(2.9)]{AL}.

\paragraph{Step 4: general profiles.}
Let $\mu$ be an arbitrary admissible profile. Following
\cite[Section~2.6]{AL}, we approximate $\nu$ by its pushforwards $\nu_j$
onto finite meshes that contain $0$, $1$, and any prescribed finite set
of levels, with mesh size tending to zero. These pushforwards are taken under the map that sends each point to
the smallest mesh point greater than or equal to it. Let
$\mu_j=s\delta_0+\nu_j$. Then $\mu_j\Rightarrow\mu$, the cumulative
functions converge in $L^1([0,1])$, and they agree with $a$ at the mesh
points. The spatial derivatives of fixed order are bounded uniformly
over step profiles at fixed $N$. We apply the $L^1$ estimate recalled in
\cref{sec:fnd-functional} to the solutions, and the same comparison
argument to their differentiated equations. Using also the uniform bound on the derivatives and
Gr\"onwall's inequality, we find that the solutions and their first and
second spatial derivatives converge uniformly. We couple the laws for $\mu_j$ and $\mu$
through the same matrix $W$ under its Gaussian reference law and the
same Brownian motion $B$. We also use the same uniform variable to
sample the terminal spin. Since the drifts are bounded and Lipschitz, Gr\"onwall's inequality
implies that $X$, $m$, and $H$ for $\mu_j$ converge to those for $\mu$
uniformly in $q$, in probability. The conditional second moments of the
terminal spin, which solve linear backward equations with drift
$ca\nabla U$, converge in the same way, and so does $C_q$. The ratios of the tilt densities for $\mu_j$ to
the density for $\mu$ converge uniformly to one, since these ratios lie
between $e^{-2sd_j}$ and $e^{2sd_j}$, where
$d_j=(cN/2)\|a_j-a\|_{L^1}\to0$. Then the densities converge in $L^1$ of
the Gaussian reference law. Since the conditional moments and spatial derivatives used here are
bounded at fixed $N$, their convergence in probability, uniformly in the
level, implies that they also converge in every finite $L^p$ norm
needed below.
Using these convergences, we pass \eqref{eq:in-martingale}, the formula for the
derivative of $\E S_q$, its continuity in $q$, and
\eqref{eq:fnd-first-variation} to the limit \cite[Section~2.6.2]{AL}.
If $\mu$ is a minimizer, we apply the limiting first variation directly
to $\mu$; the approximating step profiles need not be minimizers. Step 3
then applies to $\mu$ without change.
\end{proof}

\section{The integrated variance bound and the overlap error}
\label{sec:fnd-variance}

At a minimizer, we use the contact identities to turn the posterior
variance bound
of \cite[Proposition~2.2]{AL} into bounds that depend only on the
cumulative function $a$.

\begin{proposition}[Integrated variance and trace bounds]
\label{prop:fnd-variance}
Let $0<\rho\le1$, and let $\mu=s\delta_0+\nu$ minimize
$\mathcal F_{t,s}$. For every deterministic $k\in L^2(\nu)$,
\begin{equation}\label{eq:in-Lvar}
 \Var\left(\int k(q)L_q\,\nu(\dd q)\right)
      \le\frac1M\int\frac{k(q)^2}{a(q)}\,\nu(\dd q),
 \qquad A\le\frac1{Ms}.
\end{equation}
\end{proposition}

The first bound is \cite[Proposition~2.2, (2.16)]{AL}, and the second
is \cite[Section~2.5, (2.18)]{AL}. We briefly recall the argument. We fix
the terminal spin and apply the
Brascamp--Lieb inequality to the law of the independent Gaussian
coordinates that generate $W$ and the field. For a fixed spin, this law
is uniformly log-concave, although the
marginal law of $W$ under $\Pp$ is not Gaussian.

\begin{proof}[Proof of \cref{prop:fnd-variance}]
We first show that, for every admissible profile and every deterministic
$k\in L^2(\nu)$,
\begin{equation}\label{eq:fnd-L-general}
 \Var\left(\int k(q)L_q\,\nu(\dd q)\right)
 \le\frac1{N^2}\int k(q)^2\,\E\Tr(H_qC_q)\,\nu(\dd q),
\end{equation}
which is \cite[(2.15)]{AL}. Let $\mu$ be a step profile. For
$\varepsilon>0$, we start the field at $X_0=\sqrt\varepsilon z_0$, where
$z_0$ is a standard Gaussian vector included in the outer tilt. We
collect the independent standard Gaussian coordinates that generate $W$,
$z_0$, and the field increments in one vector $Z$. With the terminal spin
fixed, $Z$ has a density proportional to $e^{-\Phi_\sigma(Z)}$, and the
Hessian of $\Phi_\sigma$ is
$\mathcal Q=I+\int\mathcal A_q\,\nu(\dd q)\succeq I$, where
$\mathcal A_q=\nabla_Z^2U(q,X_q;W)\succeq0$ \cite[(2.1)--(2.3)]{AL}. We
derive this density for more general laws in the proof of
\cref{lem:fnd-posterior}. We work in these Gaussian reference
coordinates, in which, at the grid points $q_\ell$ of the profile,
\[
 X_{q_\ell}=\sqrt\varepsilon z_0
   +\sum_{h=1}^{\ell}\sqrt{c(q_h-q_{h-1})}\,z_h ;
\]
the tilted law is obtained by reweighting these coordinates. In this
representation, the field depends on $z_0$ only through the term
$\sqrt\varepsilon z_0$, which gives
$\nabla_{z_0}U(q,X_q;W)=\sqrt\varepsilon m_q$ and shows that
the block of
$\mathcal A_q$ for the coordinates of $z_0$ is $\varepsilon H_q$. Let $v_q$ be the vector whose $z_0$ block is
$(\sigma-m_q)/\sqrt\varepsilon$ and whose other blocks vanish.
Differentiating $L_q$ with the spin fixed, we find
\[
 \nabla_ZL_q=\frac1N\mathcal A_qv_q,\qquad
 v_q^T\mathcal A_qv_q=(\sigma-m_q)^TH_q(\sigma-m_q).
\]
For $Y_k=\int kL_q\,\dd\nu$ and $b=\int k(q)\mathcal A_qv_q\,\nu(\dd q)$,
we have $\nabla_ZY_k=b/N$. For every vector $w$, using the Cauchy--Schwarz
inequality for the forms $\mathcal A_q$ and for the integral against
$\nu$, together with $\int\mathcal A_q\,\dd\nu=\mathcal Q-I\preceq\mathcal Q$,
we have
\[
 |w^Tb|^2\le\Bigl(\int w^T\mathcal A_qw\,\nu(\dd q)\Bigr)
       \int k(q)^2v_q^T\mathcal A_qv_q\,\nu(\dd q)
 \le(w^T\mathcal Qw)\int k(q)^2v_q^T\mathcal A_qv_q\,\nu(\dd q).
\]
The choice $w=\mathcal Q^{-1}b$ yields
$b^T\mathcal Q^{-1}b\le\int k^2v_q^T\mathcal A_qv_q\,\dd\nu$. Then the
Brascamp--Lieb inequality \cite{BrascampLieb1976}, in the form
\cite[Proposition~2.2, (2.14)]{AL}, implies that the conditional variance
of $Y_k$ given $\sigma$ is at most
$N^{-2}\E[b^T\mathcal Q^{-1}b\mid\sigma]$. By the
spin-gauge symmetry of \cite[Section~2.1]{AL}, the conditional law of
$Y_k$ given $\sigma$ does not depend on $\sigma$. Using the law of total
variance, we obtain the same bound for $\Var(Y_k)$. Conditioning on $\calG_q$,
we have
$\E[(\sigma-m_q)^TH_q(\sigma-m_q)\mid\calG_q]=\Tr(H_qC_q)$. We remove the
regularization by dominated convergence as
$\varepsilon\downarrow0$, using $|L_q|\le3/2$, $\Tr(H_qC_q)\le N^2$, and
the Gaussian domination of the outer tilt. The approximation of Step~4 in
the proof of \cref{prop:fnd-contact} then extends \eqref{eq:fnd-L-general} to all
profiles, first for continuous $k$ and then by density for
$k\in L^2(\nu)$ \cite[Section~2.6.2]{AL}.

Now suppose that $\mu$ is a minimizer. Using the fact that the trace of
the product of two positive semidefinite matrices is nonnegative,
together with
$H_q-a(q)C_q\succeq0$, we have
\[
 \Tr H_q^2=\Tr\bigl(H_q(H_q-a(q)C_q)\bigr)+a(q)\Tr(H_qC_q)
 \ge a(q)\Tr(H_qC_q),
\]
and similarly
\[
 \Tr(H_qC_q)=\Tr\bigl((H_q-a(q)C_q)C_q\bigr)+a(q)\Tr C_q^2
 \ge a(q)\Tr C_q^2 .
\]
By \eqref{eq:in-contact}, $\E\Tr H_q^2\le N/c$ at every contact level,
and $\nu$ is carried by its support. Dividing by $a(q)\ge s>0$, we find that
$\E\Tr(H_qC_q)\le N/(ca(q))$ and $\E\Tr C_q^2\le N/(ca(q)^2)$ for
$\nu$-almost every $q$. Inserting the first bound into
\eqref{eq:fnd-L-general} and using $M=Nc$, we obtain the first bound in
\eqref{eq:in-Lvar}. For the second, the quantile function
\eqref{eq:fnd-quantile} and $a(Q(v))\ge s+v$ give
\[
 A\le\frac1M\int\frac{\nu(\dd q)}{a(q)^2}
  =\frac1M\int_0^{1-s}\frac{\dd v}{a(Q(v))^2}
  \le\frac1M\int_0^{1-s}\frac{\dd v}{(s+v)^2}
  =\frac1M\Bigl(\frac1s-1\Bigr)\le\frac1{Ms}. \qedhere
\]
\end{proof}

\begin{remark}\label{rem:fnd-partial-atoms}
Since $k$ in \eqref{eq:in-Lvar} is arbitrary, the bound applies to
sub-measures of $\nu$, such as the image of Lebesgue measure on an
interval of quantiles that contains only part of an atom. If $h$ is a
Borel function with $0\le h\le1$, then applying \eqref{eq:in-Lvar} with
$kh$ in place of $k$ and using $h^2\le h$, we obtain
\[
 \Var\Bigl(\int k(q)h(q)L_q\,\nu(\dd q)\Bigr)
 \le\frac1M\int\frac{k(q)^2h(q)}{a(q)}\,\nu(\dd q).
\]
\end{remark}

We now show that the overlap error splits into the three terms of the
defect.

\begin{lemma}[Decomposition of the overlap error]\label{lem:fnd-decomp}
Let $0<\rho\le1$. For every admissible profile and every $q\in[0,1]$, we
have $\Var(L_q)=j_q+\Var(S_q)/4$. If $q$ is a contact level of a
minimizer, then
\begin{equation}\label{eq:in-decomp}
 \E(R_q-q)^2=a_q^{\rm th}+2j_q+\Var(S_q).
\end{equation}
\end{lemma}

\begin{proof}
Both identities are displayed in \cite[Section~2.5]{AL} before (2.17).
Let $\xi^1=\sigma^1-m_q$ and $\xi^2=\sigma^2-m_q$. Conditionally on
$\calG_q$, these vectors are independent and centered, with covariance
$C_q$, and
\[
 R_q=S_q+\frac1N\bigl(m_q\cdot\xi^1+m_q\cdot\xi^2+\xi^1\cdot\xi^2\bigr).
\]
The three terms in parentheses have conditional mean zero, they are
pairwise orthogonal given $\calG_q$, and their conditional second moments
are $m_q^TC_qm_q$, $m_q^TC_qm_q$, and $\Tr C_q^2$. Taking expectations, we
find that for every admissible profile and every $q$,
\begin{equation}\label{eq:fnd-RS}
 \E(R_q-S_q)^2=a_q^{\rm th}+2j_q .
\end{equation}
Since $\E[R_q-S_q\mid\calG_q]=0$, this implies that
\[
 \E(R_q-q)^2=a_q^{\rm th}+2j_q+\E(S_q-q)^2
 =a_q^{\rm th}+2j_q+\Var(S_q)+(\E S_q-q)^2 .
\]
At a contact level of a minimizer, $\E S_q=q$ by \eqref{eq:in-contact},
which proves \eqref{eq:in-decomp}. Similarly,
$L_q=S_q/2+N^{-1}(\sigma-m_q)\cdot m_q$, where the second term has
conditional mean zero and conditional variance $N^{-2}m_q^TC_qm_q$ given
$\calG_q$. By the law of total variance, we have $\Var(L_q)=j_q+\Var(S_q)/4$.
\end{proof}

\section{The optimized interpolation and the defect}
\label{sec:fnd-defect}

We next recall from \cite[Proposition~2.3]{AL} that the derivative of the
interpolating functional in $t$ is an integral of squared overlap
deviations.

\begin{proposition}[Interpolation identity]\label{prop:fnd-interpolation}
Let $N\ge1$ and $0<s\le1/2$.
\begin{enumerate}[label=(\alph*)]
\item For every admissible profile $\mu$, the map
$t\mapsto\mathcal F_{t,s}(\mu)$ is Lipschitz on $[0,1]$ and continuously
differentiable on $(0,1)$, with
\begin{equation}\label{eq:fnd-fixed-derivative}
 \partial_t\mathcal F_{t,s}(\mu)
   =-\frac{\beta^2}4\int\E(R_q-q)^2\,\nu(\dd q),
\end{equation}
where the expectation is under the tilted law of $\mu$ at $t$. The
integral in \eqref{eq:fnd-fixed-derivative} lies in $[0,4(1-s)]$.
\item The function $f_{N,s}$ is Lipschitz on $[0,1]$. At every
$t\in(0,1)$ at which $f_{N,s}$ is differentiable, every minimizer $\mu$
of $\mathcal F_{t,s}$ satisfies
$f_{N,s}'(t)=\partial_t\mathcal F_{t,s}(\mu)$.
\end{enumerate}
\end{proposition}

\begin{proof}
Part (a) is \cite[Proposition~2.3 and (2.12)]{AL}. The regularity in $t$
and the extension from step profiles to general profiles are proved in
\cite[Section~2.6.3]{AL}. We sketch the computation for a step profile,
which shows where the square comes from. The parameter $t$ enters
$\mathcal F_{t,s}$ in three places: the factor $\beta\sqrt t$ of the
matrix interaction, the coefficient $c=\beta^2(1-t)$ of the field in
\eqref{eq:fnd-vector-pde}, and the penalty. We write $B(q)=\E R_q^2$
and $\Gamma(q)=\E S_q$, and note that $\E R_q=\Gamma(q)$.

\emph{Matrix interaction.} We write
$\beta\sigma^TW\sigma/2=\sum_eg_ev_e(\sigma)$, where $e$ runs over the
pairs $i<j$, with $g_e=g_{ij}$ and $v_e(\sigma)=\beta\sigma_i\sigma_j/\sqrt N$,
and over the diagonal pairs $e=(i,i)$, with $g_e=z_i$ and
$v_e=\beta/\sqrt{2N}$. By the computation of
\cref{sec:fnd-completion}, we have
$\sum_ev_e(\sigma)v_e(\sigma')=\beta^2NR(\sigma,\sigma')^2/2$. We regard
$y_e=\sqrt t\,g_e$ as Gaussian coordinates of variance $t$. By the
Gaussian heat identity, at fixed $c$, the derivative in $t$ of the outer
logarithmic moment is
\[
 \frac1{2N}\,\E\sum_e\bigl(\partial_{y_e}^2U(0,0;W)
   +s\bigl(\partial_{y_e}U(0,0;W)\bigr)^2\bigr).
\]
For a function $b(\sigma)$ of the terminal spin, we add $yb(\sigma)$ to
the terminal exponent. Differentiating the Hopf--Cole recursion twice in
$y$ at $y=0$, we find \cite[(2.10)]{AL}
\[
 \partial_yU(0,0)=\E[b(\sigma)\mid\calG_0],\qquad
 \partial_y^2U(0,0)=\E[b(\sigma)^2\mid\calG_0]
 -\int\E\bigl[\E[b(\sigma)\mid\calG_q]^2\bigm|\calG_0\bigr]\,\mu(\dd q).
\]
Applying this with $b=v_e$, summing over $e$, and using
$\E[v_e(\sigma)\mid\calG_q]^2=\E[v_e(\sigma^1)v_e(\sigma^2)\mid\calG_q]$
and $\mu=s\delta_0+\nu$, we find that the matrix interaction contributes
\[
 \frac{\beta^2}4\Bigl(1-\int B(q)\,\mu(\dd q)+sB(0)\Bigr)
 =\frac{\beta^2}4\Bigl(1-\int B(q)\,\nu(\dd q)\Bigr)
\]
\cite[(2.11)]{AL}. The term $s(\partial_{y_e}U)^2$, which comes from
the outer moment of order $s$, cancels the prescribed root mass.

\emph{Field.} As in Step~2 of the proof of \cref{prop:fnd-contact}, the
derivative of $U$ in $c$ solves the linearized equation with source
$-(\Delta U+a|\nabla U|^2)/2$, and by the Feynman--Kac formula, we have
\[
 \partial_c\Bigl(\frac1{Ns}\log\E_{\rm ref}e^{sU(0,0;W)}\Bigr)
 =\frac12\int_0^1\bigl(h(q)+a(q)\Gamma(q)\bigr)\dd q,
 \qquad h(q)=\frac1N\E\Tr H_q .
\]
By the equation for $H_q$ in Step~1 of the proof of
\cref{prop:fnd-contact} and the derivative of $\Gamma$, we have
$h'=-a\Gamma'$. Further, $h(1)=1-\Gamma(1)$, because $H_1=C_1$ and
$\Tr C_1=N-|m_1|^2$. Integrating by parts, and using that $h$ and $\Gamma$ are continuous
and that $a$ is right-continuous with $a(1)=1$, we find
\[
 h(q)+a(q)\Gamma(q)=1-\int_{(q,1]}\Gamma(r)\,\mu(\dd r)
 \qquad(0\le q\le1).
\]
The interval is open at $q$ because $a(q)$ already contains an atom of
$\mu$ at $q$. Integrating in $q$ and using $\dd c/\dd t=-\beta^2$, we
see that the field contributes
\[
 -\frac{\beta^2}2\Bigl(1-\int q\Gamma(q)\,\nu(\dd q)\Bigr);
\]
here we may replace $\mu$ by $\nu$ because the integrand vanishes at
$q=0$.

\emph{Penalty.} By Fubini's theorem, we have
$2\int_0^1qa(q)\dd q=1-\int q^2\,\mu(\dd q)$. It follows that the
penalty contributes
\[
 \frac{\beta^2}4\Bigl(1-\int q^2\,\nu(\dd q)\Bigr).
\]

Adding the three contributions, in which the constants cancel, we
obtain
\[
 \partial_t\mathcal F_{t,s}(\mu)
 =-\frac{\beta^2}4\int\bigl(B(q)-2q\Gamma(q)+q^2\bigr)\,\nu(\dd q),
\]
which is \eqref{eq:fnd-fixed-derivative}. The integral lies in
$[0,4(1-s)]$ because $|R_q-q|\le2$ and $\nu$ has mass $1-s$.

For (b), the bound in (a) implies that the functions
$t\mapsto\mathcal F_{t,s}(\mu)$ are Lipschitz with the common constant
$\beta^2(1-s)$. Their minimum $f_{N,s}$ is then Lipschitz as well. Let
$t\in(0,1)$ be a point at which $f_{N,s}$ is differentiable, and let
$\mu$ be a minimizer at $t$. For every $h$ such that $t+h\in[0,1]$, we have
$f_{N,s}(t+h)\le\mathcal F_{t+h,s}(\mu)$ and
$f_{N,s}(t)=\mathcal F_{t,s}(\mu)$. Dividing the difference by $h>0$
and letting $h\downarrow0$, we find
$f_{N,s}'(t)\le\partial_t\mathcal F_{t,s}(\mu)$; dividing by $h<0$
gives the reverse inequality. This is \cite[(2.20)]{AL}.
\end{proof}

In particular, $t\mapsto\mathcal F_{t,s}(\mu)$ and $f_{N,s}$ are
nonincreasing. We define the \emph{defect}
\[
 \calD_N(\rho)=-\frac4\beta f_{N,s}'(1-\rho)
\]
at every $\rho\in(0,1)$ such that $f_{N,s}$ is differentiable at
$1-\rho$. Since $f_{N,s}$ is Lipschitz, Rademacher's theorem implies that this
holds for almost every $\rho$, and $\calD_N$ is a measurable function
defined almost everywhere. All pointwise statements about $\calD_N$ are
understood at these values of $\rho$. The defect depends on $N$,
$\beta$, $\rho$, and $\lambda$, and not on a choice of minimizer.

\begin{corollary}[The defect identity]\label{cor:fnd-defect}
Let $N\ge1$ and $0<\lambda\le\beta/2$, and let $s=\lambda/\beta$. At
every $\rho\in(0,1)$ at which $\calD_N(\rho)$ is defined, and for every
minimizer $\mu$ of $\mathcal F_{1-\rho,s}$,
\begin{equation}\label{eq:in-envelope}
 \calD_N(\rho)=\beta(A+2J+V),\qquad
 4e_{N,\beta}(\lambda)=\int_0^1\calD_N(\rho)\dd\rho ,
\end{equation}
where $A$, $J$, and $V$ are computed under the tilted law of $\mu$. In
particular, $\calD_N\ge0$ and $e_{N,\beta}(\lambda)\ge0$.
\end{corollary}

\begin{proof}
Let $\rho$ be as in the statement, and let $\mu$ be a minimizer at
$t=1-\rho$. By \cref{prop:fnd-interpolation}(b) and
\eqref{eq:fnd-fixed-derivative}, we have
$\calD_N(\rho)=\beta\int\E(R_q-q)^2\,\nu(\dd q)$. Recall that the measure
$\nu$ is carried by the contact levels of $\mu$,
at which \eqref{eq:in-decomp}
holds. Integrating \eqref{eq:in-decomp} against $\nu$, we obtain the
first identity. For the second, since $f_{N,s}$ is Lipschitz, we have
\[
 \int_0^1\calD_N(\rho)\dd\rho=-\frac4\beta\int_0^1f_{N,s}'(t)\dd t
 =\frac4\beta\bigl(f_{N,s}(0)-f_{N,s}(1)\bigr)
 =\frac4\beta\bigl(P_\beta(s)-\phi^c_{N,\beta}(s)\bigr),
\]
where we used \eqref{eq:fnd-endpoints} in the last step. By
\eqref{eq:fnd-error-definition}, the right side is
$4e_{N,\beta}(\lambda)$; see also \cite[Proposition~2.3, (2.21)]{AL}.
\end{proof}

The individual terms $A$, $J$, and $V$ may depend on the choice of
minimizer, but by \eqref{eq:in-envelope}, their combination $A+2J+V$
does not. In the later sections, we bound $A$, $J$, and $V$ at an
arbitrary minimizer; such bounds then hold for $\calD_N(\rho)$ at almost
every $\rho$.

Near $\rho=0$, we use a cruder bound. For every $N\ge1$, $0<s\le1/2$,
$0<\rho\le1$, and every minimizer $\mu$ of $\mathcal F_{1-\rho,s}$, we
have
\[
 \int\E(R_q-q)^2\,\nu(\dd q)\le C\min\{1,M^{-1/2},M^{-2/3}s^{-1/3}\}
\]
with a numerical constant $C$. The first two terms are
\cite[Lemma~4.1, (4.1)]{AL}, and the third comes from the proof of
\cite[Theorem~1.1]{AL}; see \cite[(4.2)]{AL}. Keeping the first two
terms and using $\beta M^{-1/2}=(N\rho)^{-1/2}$, we obtain
\begin{equation}\label{eq:fnd-moving}
 \begin{gathered}
 \int\E(R_q-q)^2\,\nu(\dd q)\le C\min\{1,M^{-1/2}\},\\
 0\le\calD_N(\rho)\le C\min\{\beta,(N\rho)^{-1/2}\}
 \qquad(0<\rho\le1),
 \end{gathered}
\end{equation}
where the second bound holds at every $\rho$ at which $\calD_N(\rho)$ is
defined.

\section{The scalar trial profile and the trial budget}
\label{sec:fnd-trial}

The scalar trial profile $\gamma^{\rm tr}_\lambda=\max\{\lambda,\gamma_\beta\}$
raises the Parisi profile $\gamma_\beta$ to the floor $\lambda$. By
\cite[Lemma~7.2]{AL}, $\gamma_\beta$ grows linearly near $0$, and the two
profiles differ only on an interval of length of order $\lambda$, on
which the first variation vanishes at $\gamma_\beta$. The cost of the
floor is of order $\lambda^5$. A short lower plateau and its
fifth-order variational cost also occur in the replica calculation
of~\cite[Appendix~B, (B1)--(B3)]{ParisiRizzo2009}.
Here $\lambda=\beta s$ specifies a mass-floor constraint; we do not
identify it with the replica parameter in that calculation.
The uniform bound used here is the rigorous estimate of
\cite[Lemma~7.2]{AL}. Precisely, by
\cite[Lemma~7.2, (7.3)]{AL}, applied with $\beta_0=\beta$, which uses
$\beta>1$, there exist $\lambda_*>0$ and $C_*<\infty$, depending only on
$\beta$, such that
\[
 \beta^{-1}\mathcal P_\beta(\gamma^{\rm tr}_\lambda/\beta)
 \le E_\beta(0)+C_*\lambda^5
 \qquad(0<\lambda\le\lambda_*).
\]
We set $L_s=\max\{1,C_*\}$ and $\lambda_0=\min\{1,\lambda_*/2,\beta/4\}$.
Recall that $E_\beta$ is nondecreasing, and note that the profile with
cumulative function $\gamma^{\rm tr}_\lambda/\beta$ has mass at least
$\lambda/\beta$ at zero. We obtain
\begin{equation}\label{eq:fnd-scalar-trial-cost}
 0\le E_\beta(\lambda)-E_\beta(0)
 \le\beta^{-1}\mathcal P_\beta(\gamma_\lambda^{\rm tr}/\beta)
                      -E_\beta(0)
 \le L_s\lambda^5\qquad(0<\lambda\le2\lambda_0).
\end{equation}
The bound $\lambda_0\le\lambda_*/2$ allows us to use
\eqref{eq:fnd-scalar-trial-cost} at $2\lambda$, and $\lambda_0\le\beta/4$
ensures that $2\lambda/\beta\le1/2$ for $\lambda\le\lambda_0$, so that
the results above apply with $s=2\lambda/\beta$. The bound
$\lambda_0\le1$ is used in \cref{sec:lz,sec:b7}.

\begin{lemma}[Trial budget]\label{lem:fnd-trial}
For every $N\ge1$, $0<\rho\le1$, and $0<\lambda\le\lambda_0$,
\begin{equation}\label{eq:fnd-trial}
 0\le F_\rho(\gamma_\lambda^{\rm tr})-\min_\gamma F_\rho(\gamma)
       \le e_{N,\beta}(\lambda)+L_s\lambda^5 .
\end{equation}
\end{lemma}

\begin{proof}
The lower bound holds because $\gamma^{\rm tr}_\lambda$ is admissible.
Let $\mu^{\rm tr}$ be the admissible profile with cumulative function
$\gamma^{\rm tr}_\lambda/\beta$. By \cref{prop:fnd-interpolation}(a),
the map $t\mapsto\mathcal F_{t,s}(\mu^{\rm tr})$ is nonincreasing, and
$\mathcal F_{0,s}(\mu^{\rm tr})=\mathcal P_\beta(\gamma^{\rm tr}_\lambda/\beta)$
by \cref{sec:fnd-functional}. Hence
$F_\rho(\gamma^{\rm tr}_\lambda)\le
\beta^{-1}\mathcal P_\beta(\gamma^{\rm tr}_\lambda/\beta)$. Recall that
the optimized value $f_{N,s}$ is also nonincreasing. Using
\eqref{eq:fnd-endpoints} and \eqref{eq:fnd-error-definition}, we find
\[
 \min_\gamma F_\rho(\gamma)=\beta^{-1}f_{N,s}(1-\rho)
 \ge\beta^{-1}f_{N,s}(1)=\beta^{-1}\phi^c_{N,\beta}(s)
 =E_\beta(\lambda)-e_{N,\beta}(\lambda).
\]
Subtracting the two bounds, we obtain
\[
 F_\rho(\gamma^{\rm tr}_\lambda)-\min_\gamma F_\rho(\gamma)
 \le e_{N,\beta}(\lambda)
 +\beta^{-1}\mathcal P_\beta(\gamma^{\rm tr}_\lambda/\beta)-E_\beta(\lambda)
 \le e_{N,\beta}(\lambda)+L_s\lambda^5,
\]
where in the last step we used \eqref{eq:fnd-scalar-trial-cost} and
$E_\beta(\lambda)\ge E_\beta(0)$. This is
\cite[Proposition~7.7, (7.44)]{AL}.
\end{proof}

We call a number $\eta\ge0$ that bounds the left side of
\eqref{eq:fnd-trial} a \emph{trial budget}. By \cref{lem:fnd-trial},
$\eta=e_{N,\beta}(\lambda)+L_s\lambda^5$ is a trial budget
simultaneously for every $\rho\in(0,1]$ and every minimizer. In
\cref{sec:lz,sec:lz-root}, we use the trial budget to control the
distance between a minimizer and the scalar trial profile, and in
\cref{sec:b7}, it determines the cutoff between the two regimes of the proof.

\section{Reduction of the mean bias}\label{sec:fnd-mean-bias}

We now reduce the mean bias to the fractional-moment error by the
midpoint argument of Step~1 in the proof of \cite[Theorem~7.8]{AL}. We
define
\[
 \psi(\lambda)=\frac1{N\lambda}
       \log\E_{\rm ref}e^{\lambda(F_{N,\beta}/\beta+D_N)}\quad(\lambda>0),
 \qquad\psi(0)=\frac{p_{N,\beta}}\beta ,
\]
so that $\psi(\lambda)=E_\beta(\lambda)-e_{N,\beta}(\lambda)$ for
$0<\lambda<\beta$, by \eqref{eq:fnd-error-definition}. As a function
of the $N(N+1)/2$ independent standard Gaussian coordinates
$(g_{ij},z_i)$, the completed free energy $F^c_{N,\beta}/\beta$ is
convex, because it is $\beta^{-1}$ times the logarithm of a sum of
exponentials of linear functions. Its partial derivatives are
$N^{-1/2}\langle\sigma_i\sigma_j\rangle$ in $g_{ij}$, where
$\langle\cdot\rangle$ is the Gibbs average of $F^c_{N,\beta}$, and
$(2N)^{-1/2}$ in $z_i$. We find that the squared norm of its gradient is
at most
\[
 \frac1N\cdot\frac{N(N-1)}2+N\cdot\frac1{2N}=\frac N2,
\]
and the function is $\sqrt{N/2}$-Lipschitz. By \cite[Lemma~5.2]{AL}, which is a
special case of Chen's convexity theorem \cite[Theorem~1]{ChenConvexity},
the function
$\lambda\mapsto\lambda^{-1}\log\E_{\rm ref}e^{\lambda F^c_{N,\beta}/\beta}=N\psi(\lambda)$
extends to a convex function on $[0,\infty)$ whose value at zero is
$\E_{\rm ref}F^c_{N,\beta}/\beta=N\psi(0)$, where we used
$\E_{\rm ref}D_N=0$. In particular, $\psi$ is convex on $[0,\infty)$.

\begin{proposition}[Mean bias]\label{prop:fnd-mean-bias}
For every $N\ge1$ and $0<\lambda\le\lambda_0$,
\begin{equation}\label{eq:fnd-mean-bias}
 0\le\frac{P_\beta-p_{N,\beta}}\beta
 \le E_\beta(0)-2E_\beta(\lambda)+E_\beta(2\lambda)
                         +2e_{N,\beta}(\lambda)
 \le 32L_s\lambda^5+2e_{N,\beta}(\lambda).
\end{equation}
\end{proposition}

\begin{proof}
The lower bound is \cite[Theorem~1.1, (1.5)]{AL}; it is also Guerra's
bound \cite{Guerra2003}. Since $\psi$ is convex on $[0,\infty)$, we have
$\psi(\lambda)\le(\psi(0)+\psi(2\lambda))/2$, that is,
$\psi(0)\ge2\psi(\lambda)-\psi(2\lambda)$. Recall that
$\psi(\lambda)=E_\beta(\lambda)-e_{N,\beta}(\lambda)$. Since
$\lambda\le\lambda_0\le\beta/4$, we have $2\lambda/\beta\le1/2$. Then
\cref{cor:fnd-defect} applies at the tilt $2\lambda$ and implies that
$e_{N,\beta}(2\lambda)\ge0$, that is, $\psi(2\lambda)\le E_\beta(2\lambda)$.
Using $E_\beta(0)=P_\beta/\beta$, we obtain
\[
 \frac{P_\beta-p_{N,\beta}}\beta=E_\beta(0)-\psi(0)
 \le E_\beta(0)-2E_\beta(\lambda)+2e_{N,\beta}(\lambda)+E_\beta(2\lambda),
\]
which is the middle inequality in \eqref{eq:fnd-mean-bias}. For the last
inequality, we use $E_\beta(\lambda)\ge E_\beta(0)$ and
\eqref{eq:fnd-scalar-trial-cost} at $2\lambda\le2\lambda_0$:
\[
 E_\beta(0)-2E_\beta(\lambda)+E_\beta(2\lambda)
 \le E_\beta(2\lambda)-E_\beta(0)\le L_s(2\lambda)^5=32L_s\lambda^5 .
 \qedhere
\]
\end{proof}

In \cref{sec:b7}, we take $\lambda=N^{-7/50}$, so that
$\lambda^5=N^{-7/10}$. \Cref{prop:fnd-mean-bias} then reduces
\cref{thm:main} to the bound $e_{N,\beta}(N^{-7/50})\le C_\beta N^{-7/10}$
for all sufficiently large $N$.
